\unnumberedChapter{The Base Station}

Take a deep breath. Over the next few chapters, we are going to put together our first major LCS
 hardware module. The previous chapters introduced the message format and protocols, the core 
 library for implementing the event system, and the underlying equipment based on the DCC 
 signal standards. We then took a closer look at the individual hardware building blocks and 
 discussed how to approach the design of node firmware in general.

Now it is time to bring all these pieces together. Here comes the first and arguably most 
important hardware module: the base station. Every layout needs some kind of base station 
that acts as the central point for layout control and signal generation.

Looking at the market, there are plenty of so-called base stations available. They typically
 support several standards and communication protocols, such as DCC, mfx, LocoNet, CAN bus, 
 S88 sensor bus, and so on. Most base stations also have the power unit directly integrated. 
 They support the configuration of locomotive and stationary decoders. In short, they offer 
 an all-in-one solution. Prices are typically in the range of a few hundred Euros. With the 
 advent of Arduino, Raspberry Pi, and other low-cost controllers, there are numerous d
 o-it-yourself solutions as well.

Just to name one example, the classic \texttt{DCC++} Arduino base station, using a motor 
shield as its power unit, provides a fully functional base station for well under a hundred
 Euros. The excellent work of the JMRI community has made it possible to use this inexpensive 
 hardware with configuration software and other utilities through a dedicated 
 \texttt{DCC++} interface. The \texttt{DCC-EX} group has taken the original \texttt{DCC++} 
 project further, extending and stabilizing it for a wider range of controllers while 
 adding new capabilities. There are many other great projects out there. The appendix 
 provides some links and pointers to this work.

\unnumberedSection{Requirements}

Our base station needs to provide a number of capabilities.

First, it must be able to assemble DCC packets and generate the corresponding DCC hardware
 signals. These tasks are split between the base station, which produces the signal content, 
 and the power unit, which drives the track.

Furthermore, the base station needs to provide a way to manage several active locomotives. 
For each active locomotive, the current state is maintained and the corresponding DCC packets 
are produced. When a new locomotive session is established, a locomotive dictionary could be 
consulted to retrieve the initial function settings and other relevant information about that 
particular locomotive.

A base station implementing locomotive session and track management should also provide a serial 
command interface for managing sessions and sending commands to individual locomotives. Although 
not strictly necessary, such an interface is extremely useful for testing and debugging. It is 
also beneficial when programming DCC decoders, as some means of transferring configuration data
 to the decoder is required.

There are excellent open-source tools available that make use of an ASCII interface to send their
 commands. One example is DecoderPro from the JMRI project. Among other protocols, it supports 
 the \texttt{DCC++ / DCC-EX} ASCII interface for sending and receiving commands. Our base 
 station will therefore implement the relevant \texttt{DCC++ / DCC-EX} commands.

Base stations are often combined with layout management systems. These provide tools for 
creating track plans, displaying graphical representations of the layout, and controlling 
locomotives, turnouts, signals, and so on. RocRail, for example, is such a management system 
and can connect to many different base stations.

Our base station should be flexible enough to accommodate these management systems as well. 
Think of a gateway approach: commands from RocRail, which come in several different forms, 
are received and translated into LCS messages. Later in the book, we will have a chapter 
dedicated to LCS gateways. A gateway, as well as a base station, is a natural place to provide
 this kind of functionality.

All configuration settings, such as the number of concurrent locomotive sessions or the current
 consumption limit for a track, should be available as attributes in the node and port structures.
  A track should be represented by a port. Consequently, we will have one port controlling the 
  MAIN track output and another controlling the PROG track output.

Finally, the base station firmware should be able to host basic LCS management functions, 
such as defining a new node. This is not strictly a base station function. However, since 
every layout needs a base station, and since it might also be used to display high-level status 
information, it is a convenient place to host some central management functionality as well.

We will discuss this subject in another chapter. For now, this chapter will focus on the core 
features of the base station itself.

\unnumberedSection{Base Station Components}

The following chapters and sections will describe the various components of our base station. 
We will start with the hardware. 
In addition to the main controller components we already know from the previous chapters, we 
need two more components.

The first is the \textbf{Power Unit}, which supplies power to the DCC track. Describing the 
power unit design also provides a good opportunity to establish some standards for connectivity 
and polarity.

The second component is a \textbf{RailCom Message Detector}. RailCom is the back-channel 
communication mechanism added to the DCC protocol. If a particular DCC decoder supports RailCom, 
the base station can receive messages transmitted by that decoder through this channel.

